\summarychapter{16}{Bending the Necks of Victors}{Mission to the Barbarians}

\begin{inbrief}
When the Germanic tribes overran the Roman West, the civilization that rose on the ruins was Christian
because the invaders gave up not their weapons but their gods. The chapter describes who the
barbarians were and how they were won: through Arian missionaries, through Irish and British monks,
through the mass conversion of the Franks under Clovis, and through the Roman mission to England and
Boniface's work in Germany.
\question{What led the Germanic conquerors to adopt the religion of the people they had defeated?}
\end{inbrief}

The tribes of the northern forests sacrificed animals and worshiped nature spirits, and a missionary
had to show that the Christian powers were stronger. In the eighth century Boniface took an axe to
the sacred oak of Thor in Germany; at the first stroke, the story says, a wind from God brought it
down, the pagans were converted, and he built a chapel to St.~Peter from the wood. A miracle here and
a victory in battle there: that is how Europe became Christian.

\ssection{The Coming of the Barbarians}
With the year 476 the Christian empire of Rome ceased to exist in the West. Real power already lay with
German mercenaries. In 475 their commander Orestes made his young son Romulus Augustulus emperor; in
476 another German, Odovacar, killed Orestes, deposed the boy and ruled in his own name. No one was
shocked. The shock had come a generation earlier with Alaric, Gaiseric and Attila.

The ``barbarians,'' so called because they spoke neither Greek nor Latin, were tribes whose homes lay in Scandinavia or close to it, among them Goths, Vandals, Burgundians, Lombards, Franks, Angles and Saxons. They lived
between herding and farming and counted wealth in cattle. Tacitus describes heavy drinkers and
gamblers but praises their courage, respect for women and freedom from Roman vices; their gods survive
in our Wednesday and Thursday. They lived to fight, each warlord surrounded by warriors bound by a
personal oath, for whom outliving a fallen chief was lifelong disgrace. For centuries Romans and
Germans traded and fought across the Rhine and Danube. Germans entered the empire as slaves, settlers
and soldiers, and by the late fourth century the western army was almost wholly German.

Mass invasion began with the Huns, who crossed the Volga and subdued the Ostrogoths. The frightened
Visigoths were admitted across the Danube as allies in 376, were mistreated by corrupt officials and
rose in revolt; in 378 the emperor Valens lost his army and his life at Adrianople, a battle that
ended the legend of the invincible legions. Theodosius I held them for a time, but after his death in
395 they moved under Alaric into Italy and sacked Rome in 410. To defend Italy the Romans withdrew
their troops from the Rhine in 406 and from Britain in 407, and the tribes poured in: the Vandals
through Gaul and Spain to Africa, from where they sacked Rome in 455; the Burgundians into the Rhone
valley; the Franks across northern Gaul; the Angles, Saxons and Jutes into Britain. Only the Frankish
and Anglo-Saxon kingdoms lasted more than a few centuries. The invasions were less catastrophic than
once thought. Most of the land taken was vacant or imperial, the Germans were a minority, and in time
they lost their customs, religion and speech, which is why Italy, France and Spain keep hardly a
trace of Germanic language.

\ssection{Missions to the North}
Nominal conversion was not hard, for the tribes wanted a share in the grandeur of Rome and saw
Christianity as the Roman religion. To tame, educate and teach them, and to instill even a little
Christian behavior, was another matter. They came to orthodoxy by two roads, directly from paganism or
through Arianism. The missionary Ulfilas worked among the Visigoths for forty years, half a century
after Nicaea, and translated the Bible into Gothic---except the books of Kings, too warlike for so
fierce a people. From the Goths Arianism spread to other tribes. It was more a matter of church
structure than theology: Christ was a glorified warlord, the churches belonged to the clan, and there
was no center like Rome, whose authority these peoples resisted even after turning orthodox.

Some of the first northern Christians lived outside the empire. The Celtic Irish trace their
conversion to Patrick, early in the fifth century, who was not Irish. Raised in a Christian family in
Britain, he was carried off by Irish raiders, came to living faith as a slave herding swine, escaped
after six years, and returned as a missionary after a dream in which the children of Ireland called
him back. A century later the Irish church was mainly monastic, a form that suited a rural Celtic
society, and it became a base for mission. Columba founded the monastery of Iona off Scotland, which
gave new life to preaching in Britain; Columbanus and others went on to the Continent and founded
houses such as St.~Gall and Bobbio in Germany, Switzerland and northern Italy. The independent spirit
of these Celtic monks later troubled Roman Christianity.

\ssection{The Gospel Among the Franks}
Roman Catholic influence began among the Franks of northern Gaul, the tribe destined to shape
Christian Europe. Their founder Clovis (481--511) was the first important barbarian chief to accept
orthodox Christianity. His Burgundian wife Clothilda urged the faith on him; their first son died in
his baptismal robes, a second recovered after her prayers. Facing defeat by the Alamanni, Clovis
vowed baptism in return for victory, and won. The bishop of Rheims told him, ``Bend your neck. Worship
what you burned and burn what you worshiped.'' Three thousand of his warriors were baptized with him.

This first of many medieval mass conversions raises the question of the paganizing of Christianity.
Individual conversion after instruction, the method of nineteenth-century Protestant missions,
uproots converts from their culture. Mass conversion, which won Europe, leaves people in their
culture but brings their superstitions and conduct into the church. For Clovis Jesus was a tribal
war-god. The Franks admired Peter for drawing his sword, and military saints flourished, such as
St.~George for England and St.~James for Spain. The saints, each with a specialty---pigs, hens,
toothache, fever, sore throats---may have meant more to people than Christ, and tales of their
relics abounded, like that of the lame and the blind beggar who fled the relics of St.~Martin for
fear of being cured and losing their alms.

\ssection{Passage to England}
Shallow as it was, the Franks' conversion opened a way to Britain, where the hatred between Britons
and Anglo-Saxon invaders ruled out any mission from the Britons. The initiative was Rome's. Pope Gregory the Great, in 596, sent Benedictine monks under Augustine, later called Augustine of
Canterbury, to Kent. Queen Bertha, a Christian, obtained a hearing from King Ethelbert---outdoors,
for fear of the stranger's magic---and the king granted land at Canterbury, ever since the seat of
England's religious leader. Gregory made Augustine the first archbishop. As the Roman mission moved
north it met the Celtic mission moving south. King Oswy kept the Celtic date of Easter while his
queen kept the Roman one. At the Synod of Whitby in 664, hearing that Peter holds the keys of heaven,
Oswy chose not to offend the doorkeeper and adopted the Roman practice. From then on the British Isles
moved into the orbit of Rome.

Rooted in England, Christianity returned in strength to the Continent. Winfrid of Devonshire, known
as Boniface, received from Pope Gregory~II in 729 the task of evangelizing Germany. He converted great
numbers, brought the Irish and British monks and their converts into closer relation with Rome, and
left behind a church firmly bound to the central authority. As archbishop of Mainz he could have ended
in peace, but he went back to the still pagan Frisians and died a martyr, probably in 754, when the
Frankish prince Charles, the future Charlemagne, was twelve.

\begin{keyfacts}
\keypeople{Boniface (Winfrid); Odovacar; Romulus Augustulus; Alaric; the emperor Valens; Ulfilas;
Patrick; Columba; Columbanus; Clovis and Clothilda; Gregory the Great; Augustine of Canterbury;
Ethelbert and Bertha of Kent; King Oswy.}
\keydates{%
  \kd{376}{The Visigoths cross the Danube into the empire}
  \kd{378}{Battle of Adrianople: the emperor Valens killed}
  \kd{406--407}{Roman troops withdrawn from the Rhine and from Britain}
  \kd{476}{Odovacar deposes Romulus Augustulus: end of the empire in the West}
  \kd{481--511}{Clovis, king of the Franks, the first orthodox barbarian king}
  \kd{596}{Gregory the Great sends Augustine to England}
  \kd{664}{Synod of Whitby: England adopts Roman practice}
  \kd{729}{Boniface commissioned to evangelize Germany}
  \kd{754}{Martyrdom of Boniface in Frisia}}
\keyterms{Barbarians; Visigoths, Vandals, Franks, Anglo-Saxons; Gothic Bible; Arian Christianity;
Celtic monasticism; Iona; mass conversion; cult of the saints; archbishop of Canterbury.}
\end{keyfacts}
